\section{Анализ предметной области}

\subsection{Введение в предметную область}

Механические клавиатуры -- сегмент периферийных устройств с высокой кастомизацией. Целевая аудитория: геймеры, программисты, энтузиасты DIY. Отличаются тактильной отдачей, долговечностью и возможностью замены компонентов.

Сплит-клавиатура (split keyboard) -- это устройство ввода, разделённое на две (реже-более) отдельные части. В отличие от традиционной клавиатуры, где все клавиши расположены в едином блоке, здесь левая и правая половины независимы.

\subsection{Анализ научного материала}

Статья "<Virtual human hand model and kinematics"> посвящена моделированию человеческой руки с точки зрения биомеханики и математики. В ней показано, что рука -- это очень сложная система с большим количеством степеней свободы, где каждый палец имеет несколько суставов и может двигаться только в определённых пределах. Движения руки не произвольны: они ограничены анатомией, углами суставов и связками. При этом сама форма руки не прямая -- пальцы расположены под разными углами, а ладонь имеет естественный изгиб.

Авторы создают математическую модель руки, чтобы лучше понимать, как человек взаимодействует с объектами, и использовать это для проектирования различных устройств. Главная идея заключается в том, что рука одновременно очень подвижная и очень ограниченная: она может выполнять сложные движения, но только в рамках своей естественной биомеханики.

С этим напрямую связаны раздельные клавиатуры. Обычная клавиатура предполагает, что руки лежат прямо и параллельно, хотя в реальности анатомия руки устроена иначе. Такое положение не соответствует естественным углам пальцев и запястий, из-за чего возникает лишнее напряжение. Раздельная клавиатура, наоборот, позволяет расположить руки под более естественным углом, ближе к тому, как они устроены в реальности.

Кроме того, поскольку у суставов есть чёткие ограничения по движению, неудобная клавиатура может заставлять человека выходить за комфортные диапазоны, что со временем приводит к усталости и нагрузке. Раздельная клавиатура помогает держать движения ближе к естественным, снижая это напряжение. Также важно, что пальцы работают не полностью независимо, а как связанная система, и более естественное расположение клавиш упрощает координацию движений.

В итоге из статьи можно сделать общий вывод: устройства должны проектироваться с учётом реальной биомеханики руки. Раздельные клавиатуры как раз следуют этому принципу -- они не заставляют человека подстраиваться под неудобный инструмент, а наоборот адаптируются под форму и движения руки. Поэтому их можно считать более эргономичным и потенциально более полезным решением для длительной работы.

\subsection{Основные цели использования}

\underline{Улучшение эргономики и забота о здоровье} -- ключевые преимущества сплит клавиатур, которые позволяют расположить руки в естественном положении и снизить нагрузку на запястья и плечи при длительной работе за компьютером.

Главная задача сплит клавиатур -- снизить нагрузку на опорно-двигательный аппарат при длительной работе за компьютером.

Проблемы, которые решают:

\begin{enumerate}
\item Синдром запястного канала (карпальный туннельный синдром) -- сдавление срединного нерва в запястье из-за неестественного положения руки.
\item Тендинит -- воспаление сухожилий из-за повторяющихся движений.
\item Общие боли в запястьях, плечах и спине -- следствие неправильной позы.
\end{enumerate}

Как это работает:

\begin{enumerate}
\item Руки располагаются шире, примерно на уровне плеч -- это естественная анатомическая позиция.
\item Запястья остаются прямыми, без изгиба влево/вправо.
\item Плечи не напрягаются, локти не оттопыриваются.
\end{enumerate}

\underline{Индивидуальная настройка под пользователя} -- одно из главных преимуществ сплит клавиатур: можно регулировать расстояние между половинками и угол наклона каждой части, добиваясь максимально комфортного положения рук.

Сплит клавиатуры позволяют адаптировать рабочее место под личные особенности:

\begin{enumerate}
\item Ширина плеч: если у вас широкие плечи, раздвиньте половинки подальше. Если узкие -- сдвиньте ближе.
\item Угол наклона: многие модели дают регулировать наклон каждой половины отдельно (обычно от до ), создавая «отрицательный наклон» -- когда передняя часть ниже задней. Это снижает изгиб запястья.
\item Асимметричное расположение: если одна рука чуть длиннее или есть особенности осанки, можно настроить половинки по-разному.
\end{enumerate}

\underline{Повышение продуктивности и комфорта} при работе за компьютером достигается благодаря сплит-клавиатурам: их эргономичная конструкция снижает усталость рук, позволяя дольше сохранять высокую скорость набора текста.

\begin{enumerate}
\item Меньше усталости: благодаря естественному положению рук можно работать дольше без дискомфорта.
\item Точность набора: некоторые модели учат «чистому» слепому методу печати -- пальцы не «заходят» на чужую половину, что тренирует мышечную память.
\item Снижение ошибок: меньше случайных нажатий из-за усталости.
\end{enumerate}

\underline{Улучшение осанки} достигается благодаря использованию сплит клавиатур: их конструкция позволяет расположить руки в естественном положении, что способствует правильному положению плеч, спины и шеи при работе за компьютером.

Правильное положение рук влияет на всё тело: 

\begin{itemize}
\item плечи не поднимаются к ушам;
\item спина остаётся прямой;
\item шея не вытягивается вперёд.
\end{itemize}

Это снижает нагрузку на позвоночник и мышцы спины.

\underline{Специализированные задачи и кастомизация} становятся доступными благодаря сплит клавиатурам. Возможность индивидуальной настройки раскладки, назначения макросов и создания многоуровневых слоёв позволяет оптимизировать клавиатуру под уникальные потребности пользователя -- будь то написание кода, работа с таблицами или управление игровыми командами.

Некоторые модели (например, ErgoDox EZ) предлагают:

\begin{itemize}
\item полную перенастройку раскладки клавиш;
\item программируемые слои (несколько функций на одной клавише);
\item механические переключатели для тактильной обратной связи;
\item подсветку и макросы.
\end{itemize}

\subsection{Кому особенно полезны сплит-клавиатуры?}

Сплит-клавиатуры особенно полезны:

\begin{enumerate}
\item Офисным работникам -- тем, кто ежедневно набирает много текста.
\item Программистам -- проводят за клавиатурой по 8--10 часов.
\item Дизайнерам и редакторам -- работают с текстом и интерфейсами.
\item Геймерам -- для долгих сессий (особенно в стратегиях или MMO).
\item Людям с болями в запястьях/плечах -- профилактика или облегчение симптомов.
\item Пользователям с анатомическими особенностями -- асимметрия тела, широкие/узкие плечи.
\item Сторонникам эргономичного рабочего места -- тем, кто ценит индивидуальный комфорт. 
\end{enumerate}

\subsection{Основные модели сплит-клавиатур}

\underline{Charybdis} -- есть основные варианты: 4х6, 3х6, 3х5. На рисунке \ref{Charybdis4x6:image} представлен один из вариантов клавиатуры Charybdis 4х6.
\begin{figure}[H]
	\centering
	\includegraphics[width=0.7\linewidth]{Charybdis4x6}
	\caption{Клавиатура Charybdis 4x6 Mk2}
	\label{Charybdis4x6:image}
\end{figure}
\underline{Scylla} -- та же сама Charybdis, только без трекбола, также представлена в трёх вариантах: 4х6, 3х6, 3х5. На рисунке \ref{Scylla4x6:image} представлен один из вариантов клавиатуры Scylla в формфакторе 4х6.
\begin{figure}[H]
	\centering
	\includegraphics[width=0.7\linewidth]{Scylla4x6}
	\caption{Клавиатура Scylla 4х6}
	\label{Scylla4x6:image}
\end{figure}
\underline{Corne} -- одна из самых популярных моделей сплит-клавиатуры, ввиду чего больше кастомизации на неё, соответственно и гайдов по её сборке и настройке. На рисунке \ref{Corne:image} показана одна из вариаций клавиатуры Corne.
\begin{figure}[H]
	\centering
	\includegraphics[width=0.7\linewidth]{Corne}
	\caption{Клавиатура Corne}
	\label{Corne:image}
\end{figure}
\underline{Lily58 pro} -- клавиатура с самым большим клавиш, идеально подходит под первую сплит-клавиатуру, чтоб в дальнейшем перейти на что-то более «удобное», хотя понятие удобство у каждого своё. На рисунке \ref{Lily58pro:image} представлена один из множества вариантов данного устройства ввода.
\begin{figure}[H]
	\centering
	\includegraphics[width=0.7\linewidth]{Lily58pro}
	\caption{Клавиатура Lily58 pro}
	\label{Lily58pro:image}
\end{figure}
\underline{PNCATEHO} -- клавиатура с 10 клавишами для одной руки, способная вместить все 100\% клавиш с обычной клавиатуры. Крайне сложна в освоении, но вполне возможна. На рисунке \ref{PNCATEHO:image} показана одна из её вариаций под левую руку.
\begin{figure}[H]
	\centering
	\includegraphics[width=0.7\linewidth]{PNCATEHO}
	\caption{Клавиатура Lily58 pro}
	\label{PNCATEHO:image}
\end{figure}

\subsection{Основные варианты прошивок. Преимущества и недостатки}

\underline{QMK (Quantum Mechanical Keyboard)} -- самая популярная прошивка с открытым исходным кодом для механических клавиатур. Чаще всего используется на проводных моделях клавиатур в целом.

Ключевые возможности:

\begin{itemize}
\item полная кастомизация раскладки;
\item поддержка множества слоёв (до 32);
\item функции Hold Tap (клавиша выполняет разные действия при нажатии и удержании);
\item комбо клавиши (комбинации нескольких клавиш дают новый результат);
\item макросы;
\item RGB подсветка с эффектами;
\item OSD (экранное меню);
\item поддержка сложных модификаторов.
\end{itemize}

Совместимость:

\begin{enumerate}
\item Проводные контроллеры: Pro Micro, Elite C, Elite Pi, Liatris и аналоги.
\item Подходит для большинства самодельных и готовых сплит клавиатур.
\end{enumerate}

Способы настройки:

\begin{itemize}
\item локальная сборка через QMK Toolbox;
\item онлайн конфигуратор QMK Configurator;
\item редактирование кода вручную.
\end{itemize}

Плюсы:

\begin{itemize}
\item огромное сообщество и документация;
\item регулярные обновления;
\item гибкость настройки.
\end{itemize}

Минусы:

\begin{itemize}
\item сложность для новичков;
\item требует базовых навыков программирования для полной кастомизации.
\end{itemize}

\underline{Vial} -- графический интерфейс поверх QMK, упрощающий настройку клавиатуры. Используется как в проводных, так и в беспроводных клавиатурах.

Ключевые особенности:

\begin{itemize}
\item визуальное редактирование раскладки без программирования;
\item динамические слои (изменение раскладки «на лету»);
\item поддержка всех функций QMK;
\item кроссплатформенность (Windows, macOS, Linux);
\item синхронизация настроек между устройствами.
\end{itemize}

Совместимость: те же контроллеры, что и QMK (Pro Micro, Elite series и др.).

Плюсы:

\begin{itemize}
\item интуитивно понятный интерфейс;
\item не требует навыков программирования;
\item мгновенное применение изменений;
\item возможность «откатить» настройки.
\end{itemize}

Минусы:

\begin{itemize}
\item меньшая гибкость, чем при прямой работе с QMK;
\item зависимость от приложения Vial. 
\end{itemize}

\underline{ZMK (Zephyr Mechanical Keyboard)} -- прошивка, оптимизированная для беспроводных клавиатур на базе чипов Nordic nRF52 (например, nice!nano).

Ключевые возможности:

\begin{itemize}
\item нативная поддержка Bluetooth LE;
\item низкое энергопотребление;
\item модульная архитектура;
\item поддержка беспроводного обновления прошивки (OTA);
\item продвинутая обработка ввода с низкой задержкой;
\item встроенная поддержка поворотных энкодеров.
\end{itemize}

Совместимость:

\begin{itemize}
\item беспроводные контроллеры: nice!nano, XIAO nRF52840 и аналоги;
\item клавиатуры с Bluetooth модулями.
\end{itemize}

Плюсы:

\begin{itemize}
\item оптимальна для беспроводных сборок;
\item энергоэффективность;
\item активное развитие проекта.
\end{itemize}

Минусы:

\begin{itemize}
\item меньше документации, чем у QMK;
\item сложнее в освоении;
\item ограниченная поддержка проводных клавиатур.
\end{itemize}

\subsection{Краткая история сплит-клавиатур}

История сплит-клавиатур начинается с критики прямой раскладки, запатентованной Христофором Латамом Шолесом в 1878 году (QWERTY раскладка). Несмотря на широкое распространение, со временем она стала ассоциироваться с повышенным риском развития мышечных напряжений и болей в руках и плечах у активных пользователей.

Первое серьёзное предложение по изменению дизайна клавиатуры для улучшения эргономики было сделано в 1920-х годах, но исследования в этом направлении активизировались лишь с работами Карла Кроемера в 1960-х годах. Кроемер экспериментировал с различными углами и расположениями клавиш, демонстрируя, что модификации дизайна могут значительно снизить дискомфорт и улучшить производительность печати.

В конце XX века, на фоне роста популярности персональных компьютеров, возникла потребность в решениях для предотвращения проблем с опорно-двигательным аппаратом среди пользователей.

Одним из первых коммерческих решений стали клавиатуры Maltron, представленные в 1980-х годах. Они предложили эргономичную раскладку и сплит-дизайн для снижения нагрузки на руки.

В 1990-х годах крупные технологические компании, такие как Apple и Microsoft, разработали собственные эргономичные модели.

Apple Adjustable Keyboard -- одна из первых сплит-клавиатур, предлагавшая пользователю возможность настройки угла разделения. На рисунке \ref{AppleAdjustableKeyboard:image} она и представлена.
\begin{figure}[H]
	\centering
	\includegraphics[width=0.75\linewidth]{AppleAdjustableKeyboard}
	\caption{Клавиатура Apple Adjustable Keyboard}
	\label{AppleAdjustableKeyboard:image}
\end{figure}

Microsoft Natural Keyboard, представленная в 1994 году, стала символом эргономического дизайна клавиатур. На рисунке \ref{MicrosoftNaturalKeyboard:image} её референс и представлен.
\begin{figure}[H]
	\centering
	\includegraphics[width=0.75\linewidth]{MicrosoftNaturalKeyboard}
	\caption{Клавиатура Microsoft Natural Keyboard}
	\label{MicrosoftNaturalKeyboard:image}
\end{figure}

Исследования конца 1990-х и начала 2000-х годов подтвердили, что сплит-клавиатуры могут значительно уменьшать напряжение в руках, улучшать осанку и снижать болевые ощущения у пользователей. Эти открытия привели к тому, что к концу 2000-х годов сплит-клавиатуры стали лидерами продаж среди всех клавиатур в розничной торговле США, подтверждая их эргономическую ценность и популярность среди широкой аудитории пользователей.

На текущем этапе рынок сплит-клавиатур продолжает расти, обогащаясь новыми технологиями и инновациями. Производители внедряют беспроводные технологии, улучшенную эргономику, настраиваемую подсветку и возможность программирования клавиш, что делает сплит-клавиатуры ещё более адаптивными к индивидуальным потребностям пользователей.

Среди современных тенденций -- добавление дополнительных функций, таких как интегрированные трекпады, трекболы или энкодеры для управления объёмом звука или прокруткой.

\subsection{Строение сплит-клавиатуры на примере Charybdis 4x6 Mk2}

В данной дипломной работе остановлюсь поподробнее на клавиатуре Charybdis 4x6 Mk2, ибо под неё и будет написано программное обеспечение по ходу самой работы.

Charybdis 4x6 Mk2 -- это эргономичная сплит-клавиатура с интегрированным трекболом, предназначенная для снижения нагрузки на руки и повышения эффективности работы.

\underline{Отличается от аналогов:}

\begin{itemize}
\item модульной сборкой (возможность комбинировать половинки);
\item поддержкой hot swap;
\item открытой прошивкой (QMK/ZMK).
\end{itemize}

\underline{Актуальность темы:}

\begin{itemize}
\item рост спроса на компактные клавиатуры;
\item сложность настройки существующего ПО для новичков;
\item потребность в интуитивном интерфейсе для кастомизации;
\item отсутствие единого решения, объединяющего все функции.
\end{itemize}

Она состоит из двух отдельных половин, что позволяет регулировать угол наклона и расстояние между ними в соответствии с анатомией пользователя.

\subsubsection{Корпус}

Изготовлен методом 3D-печати из материала Prusament Galaxy Black PLA. Корпус имеет вогнутую форму (concave keywell), которая уменьшает ход пальцев и обеспечивает оптимальное расположение клавиш под руками. Доступны варианты с дополнительным тентом (наклонной панелью) под углом 30 градусов. На рисунке \ref{CaseAccessories:image} изображён корпус клавиатуры и соответствующая фурнитура для неё же.
\begin{figure}[H]
	\centering
	\includegraphics[width=0.6\linewidth]{CaseAccessories}
	\caption{Корпус с некоторыми аксессуарами}
	\label{CaseAccessories:image}
\end{figure}

\subsubsection{Печатные платы (PCB)}

Используются гибкие печатные платы толщиной 0,6 мм. В версии Mk2 применяются пользовательские платы на базе контроллера RP2040 с 16 МБ памяти. Платы совместимы с предыдущими версиями клавиатуры (MK1). На рисунке \ref{Circuit:image} они изображены.
\begin{figure}[H]
	\centering
	\includegraphics[width=0.75\linewidth]{Circuit}
	\caption{Гибкие печатные платы}
	\label{Circuit:image}
\end{figure}

\subsubsection{Контроллер}

В Charybdis 4x6 Mk2 используется микроконтроллер RP2040. Он обеспечивает обработку ввода и управление функциями клавиатуры. На рисунке \ref{RP2040:image} он и изображён.
\begin{figure}[H]
	\centering
	\includegraphics[width=0.6\linewidth]{RP2040}
	\caption{Микроконтроллер RP2040}
	\label{RP2040:image}
\end{figure}

\subsubsection{Трекбол}

Расположен под большим пальцем, что позволяет управлять курсором мыши без необходимости перемещения руки. Трекбол поддерживает левую и правую кнопку клика, прокрутку и точное перемещение. Также он состоит из шара для трекбола, корзины и печатной платы с вырезом под сенсор для мыши. На рисунке \ref{trackballbin:image} изображена корзина, на рисунке \ref{trackball:image} -- шар для трекбола стандартного размера для данной корзины.
\begin{figure}[H]
	\centering
	\includegraphics[width=0.25\linewidth]{trackballbin}
	\caption{Корзина для трекбола}
	\label{trackballbin:image}
\end{figure}

\begin{figure}[H]
	\centering
	\includegraphics[width=0.25\linewidth]{trackball}
	\caption{Шар трекбола}
	\label{trackball:image}
\end{figure}

\subsubsection{Переключатели (Свитчи)}

Клавиатура совместима с переключателями формата MX (Cherry, Gateron, Kailh и др.). Поддерживается горячая замена переключателей. На рисунке \ref{OutemuSilentPeachv3RGB:image} изображены свитчи Outemu Silent Peach v3 RGB.
\begin{figure}[H]
	\centering
	\includegraphics[width=0.5\linewidth]{OutemuSilentPeachv3RGB}
	\caption{Outemu Silent Peach v3 RGB}
	\label{OutemuSilentPeachv3RGB:image}
\end{figure}

\subsubsection{Колпачки клавиш (Кейкапы)}

Совместимы с профилями MT3, DSA, XDA, SA, PBT. Опционально можно использовать колпачки MDA. На рисунке \ref{PBTXDACherryMX:image} изображены колпачки клавиш формата PBT XDA без гравировки для Cherry MX.
\begin{figure}[H]
	\centering
	\includegraphics[width=0.5\linewidth]{PBTXDACherryMX}
	\caption{Колпачки клавиш формата PBT XDA без гравировки для Cherry MX}
	\label{PBTXDACherryMX:image}
\end{figure}

\subsubsection{Подключение}

Клавиатура может подключаться через кабель TRS и USB-C. Доступны проводная и беспроводная версии. На рисунке \ref{typec2typec:image} изображён кабель для подключения Type-c на Type-c.
\begin{figure}[H]
	\centering
	\includegraphics[width=0.5\linewidth]{typec2typec}
	\caption{Кабель Type-c на Type-c}
	\label{typec2typec:image}
\end{figure}

\subsection{Прошивка}

Используется прошивка QMK, которая позволяет гибко настраивать раскладку, макросы и другие функции. Для упрощённой настройки можно использовать онлайн-интерфейс VIA.

\subsection{Сборка Charybdis 4x6 Mk2}

Charybdis 4x6 Mk2 представляет собой мощную эргономичную клавиатуру с интегрированным трекболом, разработанную для максимального удобства и производительности пользователей. Она является одним из наиболее популярных решений среди энтузиастов механики и любителей эргономичных устройств благодаря своему уникальному дизайну и функциональности. На рисунке \ref{Charybdis 4x6 Mk2:image} представлена готовая Charybdis 4x6 Mk2.

\begin{figure}[H]
	\centering
	\includegraphics[width=1\linewidth]{Charybdis 4x6 Mk2}
	\caption{Готовая Charybdis 4x6 Mk2}
	\label{Charybdis 4x6 Mk2:image}
\end{figure}

Данная клавиатура разделена на левую и правую части соответственно. Подключение к компьютеру осуществляется по кабелю Type-c на USB, а подключение двух половинок между собой осуществляется через кабель Type-c на Type-c. На рисунке \ref{Type-c на Type-c:image} показаны оба этих кабеля.

\begin{figure}[H]
	\centering
	\includegraphics[width=1\linewidth]{Type-c на Type-c}
	\caption{кабели Type-c на USB и Type-c на Type-c}
	\label{Type-c на Type-c:image}
\end{figure}

Подключение двух половинок и подключение к компьютеру осуществляется через порты Type-c на задней части устройства. На рисунке \ref{задняя часть Charybdis 4x6 Mk2:image} они показаны. Два верхних порта отвечают за подключение двух половинок между собой, нижние два — за подключение к ПК.

\begin{figure}[H]
	\centering
	\includegraphics[width=1\linewidth]{задняя часть Charybdis 4x6 Mk2}
	\caption{задняя часть Charybdis 4x6 Mk2}
	\label{задняя часть Charybdis 4x6 Mk2:image}
\end{figure}

\subsubsection{Процесс сборки левой половинки}

Сборка левой половинки начинается с печати корпуса на 3D-принтере и установки в него электроники: платы с двумя портами Type-c, кнопкой перезагрузки, с соккетами под свитчи. На рисунках \ref{Левая половинка вид сверху:image}-\ref{Левая половинка вид снизу:image} изображена левая половинка с электроникой вид сверху и вид снизу соответственно.

\begin{figure}[H]
	\centering
	\includegraphics[width=1\linewidth]{Левая половинка вид сверху}
	\caption{Левая половинка вид сверху}
	\label{Левая половинка вид сверху:image}
\end{figure}

\begin{figure}[H]
	\centering
	\includegraphics[width=1\linewidth]{Левая половинка вид снизу}
	\caption{Левая половинка вид снизу}
	\label{Левая половинка вид снизу:image}
\end{figure}

После сборки корпус закрывается крышкой с прорезиненными ножками и наклеенной с внутренней стороны шумоизоляцией на 3 винта. На рисунке \ref{Задняя крышка вид с нижней стороны:image} изображена крышка, а на рисунке \ref{Задняя крышка вид с внутренней стороны:image} она же, но уже с наклеенной шумоизоляцией. На рисунке \ref{Задняя крышка:image} она же, но прикрученная на 3 винта к основной части.

\begin{figure}[H]
	\centering
	\includegraphics[width=1\linewidth]{Задняя крышка вид с нижней стороны}
	\caption{Задняя крышка вид с нижней стороны}
	\label{Задняя крышка вид с нижней стороны:image}
\end{figure}

\begin{figure}[H]
	\centering
	\includegraphics[width=1\linewidth]{Задняя крышка вид с внутренней стороны}
	\caption{Задняя крышка вид с внутренней стороны}
	\label{Задняя крышка вид с внутренней стороны:image}
\end{figure}

\begin{figure}[H]
	\centering
	\includegraphics[width=1\linewidth]{Задняя крышка}
	\caption{Задняя крышка}
	\label{Задняя крышка:image}
\end{figure}

После сборки основной части в сокеты под свитчи нужно проложить шумоизоляцию (свитч-пады) (рисунок \ref{Шумоизоляция (свитч-пады):image}) и следом установить сами свитчи (рисунок \ref{Свитч:image}) и энкодер (роллер) (рисунок \ref{Энкодер (роллер):image}) в правый верхний угол. После чего получится уже почти готовая к использованию половинка (рисунок \ref{Левая половинка с проложенной шумоизоляцией и свитчами:image}).

\begin{figure}[H]
	\centering
	\includegraphics[width=1\linewidth]{Шумоизоляция (свитч-пады)}
	\caption{Шумоизоляция (свитч-пады)}
	\label{Шумоизоляция (свитч-пады):image}
\end{figure}

\begin{figure}[H]
	\centering
	\includegraphics[width=0.5\linewidth]{Свитч}
	\caption{Свитч}
	\label{Свитч:image}
\end{figure}

\begin{figure}[H]
	\centering
	\includegraphics[width=0.5\linewidth]{Энкодер (роллер)}
	\caption{Энкодер (роллер)}
	\label{Энкодер (роллер):image}
\end{figure}

\begin{figure}[H]
	\centering
	\includegraphics[width=1\linewidth]{Левая половинка с проложенной шумоизоляцией и свитчами}
	\caption{Левая половинка с проложенной шумоизоляцией и свитчами}
	\label{Левая половинка с проложенной шумоизоляцией и свитчами:image}
\end{figure}

Далее нужно подготовить кейкапы (рисунок \ref{Кейкапы:image}), надев на них оринги (рисунок \ref{Оринг:image}) и всю эту конструкцию надеть на свитчи в нужной раскладке «QWERTYUIOP». Лучше всего использовать либо однотонные кейкапы, либо с двойным литьём, в противном случае краска может быстро стереться и кейкап придёт в негодность.

\begin{figure}[H]
	\centering
	\includegraphics[width=0.7\linewidth]{Кейкапы}
	\caption{Кейкапы}
	\label{Кейкапы:image}
\end{figure}

\begin{figure}[H]
	\centering
	\includegraphics[width=0.5\linewidth]{Оринг}
	\caption{Оринг}
	\label{Оринг:image}
\end{figure}

К левой половинке клавиатуры также должна прилагаться подставка для удобного положения руки при печати. Собирается она из шести деталей, мягкой подставки под запястье и семи болтов (4 с обыкновенной шляпкой и 3 с потайной) и семью гайками (4 с более тугой прокруткой и 3 с обыкновенной), после чего подставка крепится к самой половике клавиатуры (рисунок \ref{Полностью готовая левая половинка клавиатуры:image}). На рисунке \ref{Детали для сборки подставки для левой половинки клавиатуры:image} показаны все детали для сборки подставки для левой половинки клавиатуры, а на рисунке \ref{готовая подставка:image} показана готовая подставка.

\begin{figure}[H]
	\centering
	\includegraphics[width=1\linewidth]{Детали для сборки подставки для левой половинки клавиатуры}
	\caption{Детали для сборки подставки для левой половинки клавиатуры}
	\label{Детали для сборки подставки для левой половинки клавиатуры:image}
\end{figure}

\begin{figure}[H]
	\centering
	\includegraphics[width=1\linewidth]{готовая подставка}
	\caption{готовая подставка}
	\label{готовая подставка:image}
\end{figure}

\begin{figure}[H]
	\centering
	\includegraphics[width=1\linewidth]{Полностью готовая левая половинка клавиатуры}
	\caption{Полностью готовая левая половинка клавиатуры}
	\label{Полностью готовая левая половинка клавиатуры:image}
\end{figure}

\subsubsection{Процесс сборки правой половинки}

Правая половинка клавиатуры отличается не только зеркальностью, но и наличием на тамб-кластере (выведенный блок клавиш под большой палец) трекбола. Процесс сборки также начинается с печати корпуса (рисунок \ref{Корпус под правую половику:image}), пайки и установки в него электроники (рисунок \ref{Установленная электроника в правую половинку клавиатуры:image}), печатью крыши и наклейкой на неё шумоизоляции (рисунки \ref{Задняя крышка вид с нижней стороны2:image} и \ref{Задняя крышка вид с внутренней стороны2:image}) и закручиванием её на заднюю часть клавиатуры (рисунок \ref{установленная задняя крышка на правую половинку:image}). После чего также клеятся свитч-пады (рисунок \ref{Шумоизоляция (свитч-пады):image}), устанавливаются свитчи (рисунок \ref{Свитч:image}), также подготавливаются кейкапы (рисунок \ref{Кейкапы:image}) надеванием на них орингов (рисунок \ref{Оринг:image}), после чего они надеваются на свитчи. Соответственно после этого всего устанавливается шар для трекбола (рисунок \ref{Шар для трекбола:image}) и собирается подставка для удобного положения клавиатуры при работе с нею (рисунок \ref{Готовая подставка под правую половинку клавиатуры:image}) из деталей на рисунке \ref{Детали для сборки подставки под правую половинку клавиатуры:image}. После всех манипуляций получаем полноценно рабочую правую половинку клавиатуры Charybdis 4x6 Mk2 (рисунок \ref{Готовая правая половинка клавиатуры Charybdis 4x6 Mk2:image}).

\begin{figure}[H]
	\centering
	\includegraphics[width=1\linewidth]{Корпус под правую половику}
	\caption{Корпус под правую половику}
	\label{Корпус под правую половику:image}
\end{figure}

\begin{figure}[H]
	\centering
	\includegraphics[width=1\linewidth]{Установленная электроника в правую половинку клавиатуры}
	\caption{Установленная электроника в правую половинку клавиатуры}
	\label{Установленная электроника в правую половинку клавиатуры:image}
\end{figure}

\begin{figure}[H]
	\centering
	\includegraphics[width=1\linewidth]{Задняя крышка вид с нижней стороны2}
	\caption{Задняя крышка вид с нижней стороны (не прикреплённая)}
	\label{Задняя крышка вид с нижней стороны2:image}
\end{figure}

\begin{figure}[H]
	\centering
	\includegraphics[width=1\linewidth]{Задняя крышка вид с внутренней стороны2}
	\caption{Задняя крышка вид с внутренней стороны (не прикреплённая)}
	\label{Задняя крышка вид с внутренней стороны2:image}
\end{figure}

\begin{figure}[H]
	\centering
	\includegraphics[width=1\linewidth]{установленная задняя крышка на правую половинку}
	\caption{установленная задняя крышка на правую половинку}
	\label{установленная задняя крышка на правую половинку:image}
\end{figure}

\begin{figure}[H]
	\centering
	\includegraphics[width=0.5\linewidth]{Шар для трекбола}
	\caption{Шар для трекбола}
	\label{Шар для трекбола:image}
\end{figure}

\begin{figure}[H]
	\centering
	\includegraphics[width=1\linewidth]{Детали для сборки подставки под правую половинку клавиатуры}
	\caption{Детали для сборки подставки под правую половинку клавиатуры}
	\label{Детали для сборки подставки под правую половинку клавиатуры:image}
\end{figure}

\begin{figure}[H]
	\centering
	\includegraphics[width=1\linewidth]{Готовая подставка под правую половинку клавиатуры}
	\caption{Готовая подставка под правую половинку клавиатуры}
	\label{Готовая подставка под правую половинку клавиатуры:image}
\end{figure}

\begin{figure}[H]
	\centering
	\includegraphics[width=1\linewidth]{Готовая правая половинка клавиатуры Charybdis 4x6 Mk2}
	\caption{Готовая правая половинка клавиатуры Charybdis 4x6 Mk2}
	\label{Готовая правая половинка клавиатуры Charybdis 4x6 Mk2:image}
\end{figure}

\subsection{Цели, задачи, границы исследования}

Цель: выявить потребности пользователей и пробелы в существующем ПО для разработки улучшенного программного обеспечения.

Задачи:

\begin{itemize}
\item изучить технические характеристики клавиатуры;
\item проанализировать текущее ПО;
\item выявить проблемы пользователей;
\item сформулировать требования к новому ПО.
\end{itemize}

Границы исследования: прошивка (QMK/ZMK), конфигураторы (VIA/VIAL), утилиты для настройки. Не рассматриваются: аппаратная сборка, выбор переключателей.

\subsection{Анализ конкретной клавиатуры Charybdis 4x6 Mk2}

\subsubsection{Технические характеристики}

Технические характеристики клавиатуры Charybdis 4x6 Mk2:

\begin{enumerate}
\item Форм фактор: модульная (две половинки).
\item Раскладка: 24 клавиши + дополнительные для соединения половин.
\item Переключатели: поддержка hot swap (Cherry MX и аналоги).
\item Подключение: USB C (проводное), опционально Bluetooth (в модификациях).
\item Микроконтроллер: RP2040.
\end{enumerate}

\subsubsection{Конструктивные особенности}

Конструктивные особенности клавиатуры Charybdis 4x6 Mk2:

\begin{enumerate}
\item Модульная конструкция позволяет использовать как две отдельные клавиатуры или одну компактную.
\item Поддержка QMK/ZMK -- возможность глубокой кастомизации прошивки.
\item Слои (layers) -- переключение раскладок через комбинации клавиш.
\end{enumerate}

\subsubsection{Целевое назначение}

Целевое назначение клавиатуры Charybdis 4x6 Mk2:

\begin{enumerate}
\item Программирование (быстрый доступ к командам).
\item Гейминг (макросы, компактные размеры).
\item Повседневное использование (эргономичность).
\end{enumerate}

\subsection{Анализ существующего ПО}

\underline{Официальные решения:}

\begin{enumerate}
\item QMK -- открытая прошивка с поддержкой сложных функций (слои, макросы, подсветка). Требует компиляции кода.
\item VIA -- графический конфигуратор для QMK. Позволяет настраивать раскладку, макросы и подсветку без программирования.
\item ZMK -- альтернатива QMK для беспроводных клавиатур. Оптимизирована под Bluetooth.
\end{enumerate}

\underline{Сторонние разработки:}

\begin{enumerate}
\item Кастомные прошивки на GitHub с предустановленными профилями для программирования/гейминга.
\item Скрипты на Python для автоматизации настройки подсветки.
\item Утилиты для диагностики подключения.
\end{enumerate}

\subsection{Текущее ПО}

\subsubsection{Функционал текущего ПО}

Функционал текущего ПО:

\begin{enumerate}
\item Настройка раскладки и слоёв (до 8 слоёв в QMK).
\item Программирование макросов (до 64 символов).
\item Управление RGB подсветкой (эффекты, зоны).
\item Сохранение профилей во внутренней памяти.
\item Обновление прошивки через QMK Toolbox.
\end{enumerate}

\subsubsection{Оценка удобства}

Оценка удобства:

\begin{enumerate}
\item QMK: мощный, но требует навыков программирования.
\item VIA: интуитивно, но ограничен функционалом.
\item Документация разрознена (GitHub, форумы).
\end{enumerate}

\subsubsection{Основные проблемы}

Основные проблемы:

\begin{enumerate}
\item Сложность компиляции QMK для новичков.
\item Ограниченный функционал VIA (нет динамической подсветки).
\item Отсутствие шаблонов раскладок для программирования.
\item Нестабильность Bluetooth в ZMK сборках.
\end{enumerate}

\subsubsection{Типовые сценарии использования}

Типовые сценарии использования:

\begin{enumerate}
\item \underline{Программист:}

\begin{itemize}
\item переключение между QWERTY и Dvorak;
\item макросы для команд IDE;
\item подсветка синтаксиса (цвета для слоёв).
\end{itemize}

\item \underline{Геймер:}

\begin{itemize}
\item отдельные профили для игр;
\item быстрые макросы для комбо;
\item отключение лишних клавиш.
\end{itemize}

\item \underline{Повседневное использование:}

\begin{itemize}
\item медиа клавиши;
\item переключение языков;
\item минималистичная подсветка.
\end{itemize}
\end{enumerate}

\subsection{Сравнительный анализ аналогов}

Сравнительный анализ аналогов представлен в таблице \ref{analoganalis:table}

\renewcommand{\arraystretch}{0.8} % уменьшение расстояний до сетки таблицы
\begin{xltabular}{\textwidth}{|p{2.25cm}|p{3.0cm}|p{3.25cm}|p{3.25cm}|p{2.75cm}|}
	\caption{Сравнительный анализ сплит-клавиатур\label{analoganalis:table}}\\
	\hline \centrow \setlength{\baselineskip}{0.7\baselineskip} Модель & \centrow \setlength{\baselineskip}{0.7\baselineskip} ПО & \centrow \setlength{\baselineskip}{0.7\baselineskip} Глубина кастомизации & \centrow \setlength{\baselineskip}{0.7\baselineskip} Удобство & \centrow \setlength{\baselineskip}{0.7\baselineskip} Сообщество \\
	\hline \centrow 1 & \centrow 2 & \centrow 3 & \centrow 4 & \centrow 5\\ \hline
	\endfirsthead
	\caption*{Продолжение таблицы \ref{analoganalis:table}}\\
	\hline \centrow 1 & \centrow 2 & \centrow 3 & \centrow 4 & \centrow 5\\ \hline
	\finishhead
	Lily58 & QMK + VIA & Высокая & Среднее & Активное \\ \hline
	Corne & QMK & Очень высокая & Низкое\newline (требуется код) & Активное \\ \hline
	Planck & QMK + TMK & Высокая & Низкое & Среднее \\ \hline
	Charybdis & QMK/ZMK + VIA & Высокая & Зависит от ПО & Активное
\end{xltabular}
\renewcommand{\arraystretch}{1.0} % восстановление сетки

\subsection{Технологический анализ}

Платформы:

\begin{enumerate}
\item QMK: C код, компиляция через Makefile. Поддержка всех функций клавиатуры.
\item ZMK: на базе Zephyr RTOS. Оптимизация для Bluetooth.
\item Альтернативы: TMK (устаревшая), кастомные прошивки на Rust.
\end{enumerate}

Инструменты:

\begin{enumerate}
\item Компиляторы: GCC, Clang.
\item Конфигураторы: VIA (Electron), QMK Configurator (веб).
\item Отладка: QMK Debugger, логирование через Serial.
\end{enumerate}

Интерфейсы:

\begin{enumerate}
\item USB HID -- стандартная передача нажатий.
\item Bluetooth LE -- для беспроводных версий (ZMK).
\item Протоколы: HID over GATT (BLE), USB descriptors.
\end{enumerate}

Тенденции:

\begin{itemize}
\item облачные конфигураторы (сохранение профилей в аккаунте);
\item ИИ ассистенты для рекомендаций раскладки;
\item интеграция с ОС (Windows/macOS/Linux).
\end{itemize}

\subsubsection{Выявление проблем и ограничений}

Проблемы ПО:

\begin{enumerate}
\item QMK: порог входа для новичков (требуется знание C);
\item VIA: нет поддержки динамических эффектов подсветки;
\item ZMK: баги в Bluetooth стеке (потеря соединения);
\item Отсутствие единой документации.
\end{enumerate}

Аппаратные ограничения:

\begin{itemize}
\item объём памяти микроконтроллера (ограничение на количество макросов);
\item задержка ввода при Bluetooth подключении.
\end{itemize}

Риски:

\begin{itemize}
\item несовместимость кастомных прошивок с обновлениями;
\item уязвимости в Bluetooth протоколах.
\end{itemize}